\documentclass[letterpaper]{article}
\usepackage{aaai}

\usepackage{times}
\usepackage{helvet}
\usepackage{courier}
\usepackage[T1]{fontenc}
\usepackage[utf8]{inputenc}
\usepackage{graphicx}
\usepackage{booktabs}
\usepackage{subcaption}
\usepackage{placeins}
\usepackage{amsmath}
\usepackage{url}

% Numbered citations (overrides the AAAI "named" citation macros loaded above).
\usepackage[numbers,sort&compress]{natbib}
\bibliographystyle{plainnat}

% Clickable cross-references: citations, figures, tables and sections.
\usepackage[dvipsnames]{xcolor}
\definecolor{linknavy}{RGB}{20,70,140}
\usepackage[colorlinks=true,
            linkcolor=linknavy,
            citecolor=linknavy,
            urlcolor=linknavy,
            breaklinks=true]{hyperref}

\raggedbottom
\frenchspacing
\setlength{\pdfpagewidth}{8.5in}
\setlength{\pdfpageheight}{11in}
\pdfinfo{
/Title (Inverse Folding Models Sample Sequence Space Shaped by MHC Anchor Identity and TCR Context)
/Author (Sergio Mares)
/Keywords (inverse folding, ProteinMPNN, ESM-IF1, LigandMPNN, pMHC, TCR, protein design, IEDB, SKEMPI, MHCflurry, ESM-Cambrian)
}
\setcounter{secnumdepth}{2}
\renewcommand{\topfraction}{0.9}
\renewcommand{\bottomfraction}{0.9}
\renewcommand{\textfraction}{0.07}
\renewcommand{\floatpagefraction}{0.7}
\renewcommand{\dblfloatpagefraction}{0.7}
\renewcommand{\dbltopfraction}{0.9}
\setcounter{topnumber}{4}
\setcounter{bottomnumber}{4}
\setcounter{totalnumber}{8}
\setcounter{dbltopnumber}{4}
\begin{document}

\title{Inverse Folding Models Sample Sequence Space Shaped by MHC Anchor Identity and TCR Context}
\author{Sergio E. Mares$^{1,2}$ \\
{\small $^1$Adimab, Mountain View, CA 94043, USA} \\
{\small $^2$Center for Computational Biology, University of California, Berkeley, CA 94720, USA}}
\nocopyright
\maketitle

\begin{abstract}
\begin{quote}
Inverse-folding models have emerged as potent tools for designing protein sequence from structure, and the degree to which they capture the biophysical constraints of a specific molecular interface stands as an open question. TCR cross-reactivity has caused fatal toxicity in the clinic, yet the peptide sequence space is far too large to screen experimentally. We therefore ask whether inverse-folding models can recover the native peptide from structure alone, and what their sampling distribution reveals about the constraints that give rise to cross-reactivity. We applied four inverse-folding models to fifty solved pMHC-TCR structures under two structural contexts, generating 3.9 million peptide designs. Every allele-matched design was scored against two independent binding predictors, and we characterized the full sampling distribution each model induces over peptide sequence space. No model recovered the native peptide. Recovery is instead bimodal across the two anchor positions that dominate MHC binding. The C-terminal anchor P$\Omega$, though structurally load-bearing, is largely missed, while the position-2 anchor recovers at or above the interior average. Removing the T-cell receptor increases sampling diversity and raises predicted binding affinity for all four models, though the size of that affinity gain differs almost sixfold between them. It also lowers recovery at experimentally validated contact positions, and leaves every model that reports a confidence score less confident in its own designs. Both effects indicate that the models read the receptor as part of the interface. This suggests that inverse-folding models encode peptide compatibility as a set of position-specific structural constraints rather than a single optimal sequence. Our investigation establishes the sampling distribution as a probe of the constraints governing peptide recognition, and highlights both the promise and the limitations of structure-based models for predicting cross-reactivity.
\end{quote}
\end{abstract}
\section{Introduction}
T-cell receptor (TCR) engineering therapies promise a transformative future for cancer treatment, and the first reached approval in 2024~\citep{fda2024tecelra,hong2023mage}, yet cross-reactivity, the recognition of unintended peptide-MHC (pMHC) complexes, remains a fundamental barrier to scaling them into personalized medicine~\citep{mason1998crossreactivity,wooldridge2012million}. An affinity-enhanced MAGE-A3 receptor that had cleared extensive preclinical testing ultimately caused fatal cardiac toxicity in two patients, cross-recognizing an unrelated pMHC derived from titin~\citep{linette2013cardiovascular,cameron2013titin}.

Experimental methods therefore sample only a fraction of the pMHC-TCR interaction space, which is combinatorially vast, roughly 
20$^L$  for a peptide of length L, yet far from uniform.
The anchor sites collapse the entropy at fixed positions while leaving the remainder variable~\citep{falk1991allele,rammensee1999syfpeithi}, so the effective dimensionality of the presented distribution falls far below its combinatorial ceiling. It is this low-dimensional structure that makes the problem learnable, yet the same data limitations that bound experimental sampling also limit the models trained on that data~\citep{odonnell2020mhcflurry,reynisson2020netmhcpan}, motivating approaches that exploit structural rather than purely sequence-based signals.

Inverse-folding models sample sequences from a distribution conditioned on structure~\citep{dauparas2022proteinmpnn,hsu2022esmif}, and on the receptor side of this interface, the approach has recent precedent. Comparing ProteinMPNN and ESM-IF1 against physics-based Rosetta design~\citep{alford2017rosetta} across 38 solved TCR-pMHC complexes, prior work~\citep{ribeirofilho2024exploring} found that both inverse-folding models recover native CDR3 interface residues well above chance, establishing that inverse folding captures real, non-trivial signals from a TCR-pMHC interface.

Here, we turn the same class of models toward the peptide itself. The point estimate asks whether the model recovers the native peptide the structure was solved with; the distribution asks which other peptides the model treats as compatible with the same groove and the same receptor, precisely the variant and off-target sequences that underlie cross-reactivity. Because inverse-folding models score structural compatibility rather than a labeled binding phenotype, that distribution provides a first-pass map of the tolerated peptide neighborhood around the native ligand, and, because MHC-groove geometry and TCR-contact geometry occupy structurally distinguishable positions along the peptide, acts as a probe of which positions the model treats as side-chain-constrained versus permissive.
We therefore ask whether inverse-folding models recover the native peptide, and whether the distribution they induce over sequence space concentrates on the tolerated neighborhood underlying cross-reactivity and off-target recognition.

\section{Materials and Methods}
\label{sec:methods}

\subsection{Panel and models}
Fifty pMHC-TCR crystal structures were retrieved from the Protein Data Bank~\citep{berman2000pdb} (Supplementary Table~\ref{tab:pdblist}). Every structure includes a bound peptide, the MHC chains, and a complete TCR $\alpha\beta$ heterodimer~\citep{bjorkman1987structure} resolved in the same asymmetric unit with peptides of various lengths (Supplementary Figure \ref{fig:supp1a}); the class I entries additionally include $\beta_2$-microglobulin where it was resolved. Forty structures are HLA-A*02:01. The remaining ten are one HLA-B*51:01 (4MJI), one HLA-B*08:01, one HLA-B*35, two mouse H-2L$^d$, and five MHC class II complexes, all of which were excluded from the allele-specific binding predictors. The average crystallographic resolution is 2.51 $\AA$ (Supplementary Figure \ref{fig:supp1b}).

The panel is not fifty independent observations. It contains 33 distinct peptides and 26 distinct CDR3$\beta$ sequences, with SLLMWITQC solved in five crystals and several receptors (A6, 1G4, DMF5) recurring, so a pooled statistic weights an epitope by how often it happens to have been crystallised. We therefore also report a non-redundant set: human MHC class I only, one structure per TCR family (CDR3$\beta$ clustered so clonotypes differing by one or two substitutions collapse together), and one structure per peptide, taking the best-resolution member of each group as its representative so the choice is independent of any score. This leaves 30 structures over 26 peptides when further restricted to HLA-A*02:01 (Supplementary Figure~\ref{fig:supp_dedup}).

The four models are ProteinMPNN~\citep{dauparas2022proteinmpnn}, a message-passing GNN encoding the fixed backbone as a $k$-nearest-neighbor graph and decoding sequence autoregressively; ProteinMPNN (no MHC), the same architecture retrained with MHC/TCR-class structures excluded; ESM-IF1~\citep{hsu2022esmif} and LigandMPNN~\citep{dauparas2025ligandmpnn} (Supplementary Table~\ref{tab:models}).

Each of the four tools samples only the peptide chain, with the MHC, $\beta_2$-microglobulin, and (where present) TCR chains held fixed as structural context. Two context conditions were run for every structure, MHC + $\beta_2$m + TCR$\alpha\beta$ and the TCR removed from the input structure entirely. Each condition was sampled to 10,000 raw designs at a fixed sampling temperature of $T=0.1$.

\section{Results}

\subsection{Models Sample Strong Position-Specific Sequence Bias}
\label{sec:sampling}

To establish how much of the peptide sequence space these models actually explore, we sampled 9,984 designs per (structure, model) cell in the pMHC+TCR context and counted how many were distinct. An average of only 79 were unique, $0.81\%$ of the draws, ranging from 48 for ESM-IF1 to 105 for ProteinMPNN (no MHC) with 89 for ProteinMPNN and 72 for LigandMPNN, so these models return redundancy rather than diversity (Supplementary Figure~\ref{fig:supp2}). Sequence logos~\citep{schneider1990logos} show every
model imposes strong position-specific restriction independent of recovery (Supplementary Figures~\ref{fig:supp3}--\ref{fig:supp4d}), a property of the sampling distribution rather than evolutionary conservation, tracing the same contacts seen in the crystals, where 29 of 82 MHC positions contacting the peptide anywhere in the panel do so in at least 40 of the 50 (Supplementary Figure~\ref{fig:supp1c}).

\subsection{The two MHC anchors impose different constraints and recover differently}
\label{sec:recovery}

To determine which peptide positions carry recoverable signal, we measured per-position recovery of each
crystal's own native residue, pooled across fifty structures and four models. Recovery is bimodal rather
than uniform, and the two positions where this matters most are the two positions immunology already
identifies as anchors~\citep{falk1991allele,madden1995three} (Figure~\ref{fig:fig1}). P2, seated in pocket B, constrains a specific side chain
and recovers well above the interior average ($0.65$ vs.\ $0.50$), whereas P$\Omega$, seated in pocket F,
imposes a permissive backbone constraint and recovers significantly \emph{below} it ($0.26$ vs.\ $0.50$,
$p=0.00021$). Both are structurally indispensable, yet only P2 gives a design model a signal it can
exploit. Neither result is an artefact of the panel's redundancy. On the non-redundant HLA-A*02:01 set (30 structures, 26 peptides, one per TCR family) P2 still recovers above interior ($0.64$ vs.\ $0.52$, $p=0.010$) and P$\Omega$ still falls below it ($0.17$, $p<10^{-4}$), and the P$\Omega$ deficit in fact deepens as repeated epitopes are removed, from $0.25$ on the full panel to $0.17$ here (Supplementary Figure~\ref{fig:supp_dedup}).

To ask whether that asymmetry reflects side-chain identity alone or a coarser physicochemical preference,
we compared each native anchor residue against its own peptide's interior positions on the
Kyte-Doolittle hydrophobicity scale~\citep{kyte1982simple}, and separately scored every design for whether it preserved the
native residue's side-chain class (Supplementary Figure~\ref{fig:supp12}). Both anchors are markedly more
hydrophobic than the interior they sit alongside ($3.38$ at P2 and $3.80$ at P$\Omega$, against an
interior mean of $0.86$; $p<0.0001$ for both, paired Wilcoxon across 40 structures), confirming
that both pockets recruit hydrophobic residues. The two anchors nonetheless part company on how tightly
that preference is enforced in the designs. P2 designs preserve the native side-chain class well above
the interior baseline ($86.1\%$ vs.\ $70.8\%$, $p=0.0002$), whereas P$\Omega$ designs do not
($67.2\%$ vs.\ $70.8\%$, $p=n.s.$), even though P$\Omega$ designs do show significantly reduced
hydrophobicity variance relative to the interior ($p=0.0032$). P$\Omega$ therefore behaves as a pocket
that biases toward a hydrophobic character without fixing which hydrophobic residue satisfies it, which
is what a low exact-identity recovery at a structurally load-bearing position should look like. 
% We note
% that a discrete side-chain classification counts cysteine as polar despite its hydrophobic
% Kyte-Doolittle value, and four of the nineteen native peptides carry cysteine at P$\Omega$, so the
% continuous scale is the more reliable of the two readouts here.

\begin{figure*}[tb]
\centering
\includegraphics[width=\textwidth]{figures/paper/fig1_per_position_recovery_9mer.png}
\caption{\textbf{Per-position recovery is bimodal, and the two MHC anchors diverge.} Recovery
of the native residue across the 32 nonamer crystals and 4 models, one panel per peptide position
(pMHC+TCR context). The violin gives the distribution shape and every (crystal, model) value is
overlaid as a point. Recovery is sharply bimodal at most positions, piling up near 0 and 1 rather
than spreading around a mean. Positions are only comparable within a peptide length, so each length
is shown separately; the remaining lengths are in Supplementary
Figures~\ref{fig:supp_len10} and~\ref{fig:supp_len13}.}
\label{fig:fig1}
\end{figure*}


Averaged across the entire peptide, the four models fall within a narrow band of mean recovery ($0.43$
to $0.47$ in full context and $0.29$ to $0.36$ without the TCR; Supplementary Figure~\ref{fig:supp5}),
narrow enough that the models might appear indistinguishable by inspection. A Friedman test indicates
otherwise in both contexts ($\chi^2=13.13$, $p=0.004$, full context; $\chi^2=19.38$, $p<0.001$, mhconly),
so this apparent similarity reflects proximity in magnitude rather than statistical equivalence. As an
aside, on 3QDG the central Gly-Ile-Gly motif at positions 4--6, well characterized in the literature~\citep{borbulevych2011tcrs},
recovers essentially perfectly here ($1.00$, $0.73$, $1.00$ vs.\ $0.50$), though we have not attempted a
systematic motif annotation across the full panel.

The relationship between resolution and recovery illustrates how pooling across a heterogeneous panel
can understate a genuine effect. Pooled across all fifty structures, the correlation is modest,
and only one model survives correction individually ($r=-0.29$ pooled, $p=0.040$; ESM-IF1 alone,
$r=-0.37$, $p=0.008$; Supplementary Figure~\ref{fig:supp6}). However, five of the fifty structures
(2P5W, 2BNR, 2F53, 2F54 and 2PYE) are independent crystals of the identical peptide (SLLMWITQC, the
NY-ESO-1 epitope~\citep{chen2005structural,sami2007high}) differing primarily in resolution
($1.9$--$2.7$\,\AA). Holding peptide identity fixed in this way removes a confound the pooled panel
cannot. Declining resolution drives recovery down for ProteinMPNN (no MHC) ($r=-0.97$, $p=0.005$)
and, less clearly, ESM-IF1 ($r=-0.84$, $p=0.075$), and drives unique design counts up for both
ProteinMPNN variants ($r=0.97$, $p=0.007$ and $r=0.95$, $p=0.012$; Supplementary
Figure~\ref{fig:supp7}). It is not universal: LigandMPNN moves in the opposite direction on both
measures ($r=+0.19$ and $r=-0.20$, both n.s.), so the pooled-panel estimate is diluted by
heterogeneity rather than simply attenuated.

\subsection{Most peptides sampled by inverse-folding models are predicted to be presented}
\label{sec:affinity}

Having established the recovery rates of the native peptide and the sequence space each structure
samples, we next asked whether those sampled peptides would actually be presented, scoring every unique
design from the forty HLA-A*02:01 structures for predicted HLA-A2 binding with two unrelated
predictors, MHCflurry~\citep{odonnell2020mhcflurry} and ESMCBA~\citep{esmcba,esmc2024}, neither of which had seen this panel (Figure~\ref{fig:fig2}).

\begin{figure}[tb]
\centering
\includegraphics[width=0.5\textwidth]{figures/paper/fig2_esmcba_vs_mhcflurry.png}
\caption{\textbf{MHCflurry and ESMCBA predictions of designs} log10(predicted
IC50, nM) with pMHC+TCR context, one panel per model. Dashed lines mark the 500\,nM; The percentage is the fraction of designs under 500\,nM by both predictors at once.}
\label{fig:fig2}
\end{figure}

The two predictors agree closely on the same designs ($r=0.83$--$0.93$ per model in full context, up to
$0.94$ without the TCR), and agreement between two methodologically unrelated predictors constitutes
stronger evidence than either predictor alone. On this shared landscape, projected jointly into one UMAP
embedding space~\citep{mcinnes2018umap} (Supplementary Figure~\ref{fig:supp11}), ProteinMPNN and LigandMPNN consistently sample
toward higher predicted affinity than ESM-IF1. At the 500\,nM cutoff~\citep{sette1994relationship}, 52--64\% of ProteinMPNN-family and
LigandMPNN designs qualify by MHCflurry, compared with 41\% for ESM-IF1 ($\chi^2=238.50$,
$p=2\times10^{-51}$; see also Section~\ref{sec:tcr}). This pattern is better characterized as ESM-IF1 being a
consistent low-affinity outlier than as a general superiority of ProteinMPNN over ESM-IF1, since the
other three tools do not share this deficit.

\begin{figure*}[tb]
\centering
\includegraphics[width=\textwidth]{figures/paper/fig3_tcr_context_benefit_heatmap.png}
\caption{\textbf{Where the TCR-context benefit actually concentrates.} Recovery(full) $-$
recovery(mhconly), per crystal per model, crystals ordered by resolution (best to worst, left to right).}
\label{fig:fig3}
\end{figure*}

\subsection{Without the TCR, inverse-folding models sample more widely but less confidently}
\label{sec:tcr}

To understand the effect of the TCR on what these models propose, we resampled every structure with the
TCR deleted from the input and compared the resulting designs against the full-complex designs above.

Every model explores substantially more sequence space without the TCR, with unique design counts roughly
doubling or tripling (ProteinMPNN 2047 $\to$ 6637; ProteinMPNN (no MHC) 2363 $\to$ 4478; ESM-IF1 1783 $\to$
3459; LigandMPNN 1635 $\to$ 4426; Supplementary Figure~\ref{fig:supp5}). This larger space skews toward
higher predicted affinity, but not uniformly across models. ProteinMPNN (no MHC) shows the largest shift
(strong-binder rate 60.3\%$\to$86.0\%), LigandMPNN a smaller but genuine one ($52.2\%\to62.6\%$), while
ProteinMPNN and ESM-IF1 show little change (Supplementary Figure~\ref{fig:supp10}). The affinity gain
associated with TCR removal thus appears concentrated in the two models already most context-sensitive,
rather than reflecting a general consequence of removing the TCR.

A related but distinct signal comes from each model's own confidence in the sequences it designs,
independent of either external predictor. ProteinMPNN and ProteinMPNN (no MHC) report an average
per-residue negative log-likelihood for each design, for which lower values indicate higher confidence;
LigandMPNN reports an overall confidence score, for which higher values indicate higher confidence.
Pooled across all forty HLA-A*02:01 structures, removing the TCR shifts this internal score for every
model, and in every case toward lower self-assessed confidence (ProteinMPNN mean $1.016\to1.295$;
ProteinMPNN (no MHC) $0.925\to1.128$; LigandMPNN $0.413\to0.321$; Mann-Whitney $p<10^{-200}$ for all
three; Supplementary Figure~\ref{fig:supp9a}). This shift is not uniform across the panel. Broken out per
structure and Bonferroni-corrected, it reaches significance in 15--17 of 19 structures depending on the
model, but not in the remaining two to four (Supplementary Figure~\ref{fig:supp9b}). That a model can
become less confident in its own designs by its own internal metric while simultaneously being scored as
a stronger binder by an external predictor that never saw those designs indicates that the two signals
are not redundant, since one reflects how comfortably a sequence fits the model's learned distribution, the
other how likely it is to bind, and removing the TCR does not move them the same way. This distinction is
more informative than amino-acid-level recovery or composition alone, since it speaks to what the model
itself, not just an external predictor, "believes" about the sequences it is generating.

\subsection{Recovery Loss Concentrates at Experimentally Validated TCR-Contact Positions}
\label{sec:skempi}

To understand how these residue interactions contribute functionally to TCR binding, we tested whether
the recovery lost without the TCR concentrates at positions that experimental data already implicate in
receptor engagement. We cross-referenced SKEMPI 2.0~\citep{jankauskaite2019skempi} and IEDB~\citep{vita2019iedb} against each structure's
crystallized TCR clone (matched by CDR3 sequence) and identified 21 (structure, position) pairs with an
experimentally measured mutation that alters TCR binding. Pooling recovery across exactly these
positions, TCR-present designs recover the native residue $54$--$58\%$ of the time, compared with
$34$--$38\%$ for TCR-absent designs (Figure~\ref{fig:fig3}), a consistent, though not universal, gap; a
minority of individual positions move in the opposite direction.

That comparison rests on 21 curated positions. The extended panel supports a more direct version of
the same test, against measured binding energies rather than curated annotations. Restricting SKEMPI
to single mutations inside the epitope itself yields 55 (complex, position) pairs across 17 complexes
with a measured $\Delta\Delta G$. At positions where mutation costs at least 1\,kcal/mol, the recovery
gain from having the TCR present is $+0.150$, against $+0.018$ elsewhere ($p=0.047$, Mann-Whitney;
Supplementary Figure~\ref{fig:supp_skempi_pooled}). With the TCR deleted the distinction disappears
outright ($0.243$ vs.\ $0.172$, $p=0.85$), which is the null this account predicts. We report the
direction rather than the effect size: critical positions lead in every subset we examined, but the
pooled $p$ rises to $0.19$--$0.22$ when the three same-epitope H-2A crystals or the class II entries
are excluded, and survives only 6 of 17 leave-one-complex-out folds. Half the positions come from a
single epitope, so this is a consistent direction on an underpowered sample rather than an
established effect.

A second, independent line of evidence separates what these scores do encode from what they do not.
6AM5 and 6AMU present the same DMF5 receptor and the same HLA-A*02:01 groove with two different
peptides, so threading one peptide set onto both backbones isolates conformation from sequence.
Across all three scoring models, each structure's own crystallized peptide ranks $2.3$ of 13 on its
own backbone and $11.0$ of 13 on the other (Figure~\ref{fig:fig4}). The score
therefore reads the backbone it is conditioned on, not the peptide sequence alone. Measured affinity
behaves differently. The same 13 peptides carry published $K_D$ and thermal-stability values, nine of
them as an actual measurement rather than a bound. Across those nine the score carries no
information about affinity on either backbone, for any of the three models
($|\rho|\leq0.32$, all $p>0.4$; Supplementary Figure~\ref{fig:supp_6am_kd}). Including the four
peptides reported only as $>500\,\mu$M or as not detected raises the apparent correlation to about
$\pm0.5$ and reverses its sign between the two backbones, but that follows from where the assay
ceiling falls relative to the two peptide families rather than from the scores, which is why those
peptides are excluded. The distinction matters for how these models should be used: they report
which peptides are compatible with a given groove and backbone, the quantity relevant to
cross-reactivity, and not a ranking of binding strength.


% Declared ahead of the Discussion so it can take the top of a column there. Placed after
% the Discussion it could only go at the bottom, which left the facing column short.
\begin{figure}[t]
\centering
\includegraphics[width=\columnwidth]{figures/paper/fig4_dmf5_two_backbones.png}
\caption{\textbf{Each backbone favours its own crystallized peptide.} The same 13 DMF5
peptides threaded onto 6AM5 (blue) and 6AMU (orange) and ordered by their score on 6AM5,
shown as the mean of three models with a $\pm$1 SD band, with rings and dotted lines marking
each structure's own peptide.}
\label{fig:fig4}
\end{figure}

\section{Discussion}
Inverse-folding models have rapidly advanced structure-based protein design~\citep{dauparas2022proteinmpnn,hsu2022esmif,dauparas2025ligandmpnn}, and their ability to recover native interface residues above chance~\citep{ribeirofilho2024exploring} invites the question of whether they capture the biophysics of an interface or merely reproduce its most constrained positions. High recovery is easily read as the former. Our results point to the latter, since knowing that MHC anchor positions matter explains little on its own. P2 and P$\Omega$ are both structurally indispensable yet behave oppositely under inverse folding, one a side-chain readout, the other a backbone constraint tolerant of many side chains. This finer distinction is what a sequence-design model can meaningfully report.
The same caution against pooling applies to resolution, where a weak correlation across fifty structures conceals a model-specific effect once peptide identity is controlled for, using the five SLLMWITQC/NY-ESO-1 crystals as a same-peptide replicate set. TCR removal resists a single verdict for the same reason. Diversity and predicted affinity both increase for every model, but by margins that differ almost sixfold between them, and recovery decreases on average without decreasing everywhere, because whether the model requires the TCR was never one question.

% \section{Limitations and scope}
% \label{sec:limitations}

% \begin{itemize}
% \item \textbf{$T{=}0.1$ only.}
% \item \textbf{Allele skew.} Nineteen of twenty structures are HLA-A*02:01; the one exception (4MJI,
% HLA-B*51:01) is also the hardest structure in the panel by recovery, and with $n=1$ we cannot separate
% allele identity from TCR identity or structure-specific difficulty.
% \item \textbf{Predicted affinity is not measured affinity.} Agreement between MHCflurry and ESMCBA
% constitutes real evidence, but neither substitutes for an experimental binding assay.
% \item \textbf{Not all TCR-binding evidence decouples MHC loading from TCR engagement equally well.} Only
% SPR-based dissociation-constant measurements are performed on an already-folded pMHC complex; functional
% assays (cytotoxicity, activation) cannot distinguish a peptide that ceased binding MHC from one that
% ceased being recognized by the TCR.
% \item \textbf{Recovery is not validated binding.} It measures agreement with a crystallized peptide, not
% experimentally confirmed activity of the designed sequences.
% \item \textbf{The MEL5/ELAGIGILTV case study is outside the 20-structure panel} and uses a separate,
% earlier design campaign; it's a sharper single-structure illustration, not panel-scale evidence.
% \end{itemize}

% \section{Data and code availability}
% \label{sec:data}


\FloatBarrier

\bibliography{references}


\appendix
\section{Supplementary Figures}

\begin{table*}[p]
\renewcommand{\thetable}{S1}
\centering
\small
\begin{tabular}{lllll}
\toprule
PDB & Peptide & Length & Allele & Resolution (\AA) \\
\midrule
2P5W & SLLMWITQC & 9 & HLA-A*02:01 & 2.20 \\
1QSF & LLFGYPVAV & 9 & HLA-A*02:01 & 2.80 \\
1QRN & LLFGYAVYV & 9 & HLA-A*02:01 & 2.80 \\
2BNR & SLLMWITQC & 9 & HLA-A*02:01 & 1.90 \\
2GJ6 & LLFGKPVYV & 9 & HLA-A*02:01 & 2.56 \\
2F53 & SLLMWITQC & 9 & HLA-A*02:01 & 2.10 \\
2F54 & SLLMWITQC & 9 & HLA-A*02:01 & 2.70 \\
3QDG & ELAGIGILTV & 10 & HLA-A*02:01 & 2.69 \\
3QEQ & AAGIGILTV & 9 & HLA-A*02:01 & 2.59 \\
3QFJ & LLFGFPVYV & 9 & HLA-A*02:01 & 2.29 \\
3GSN & NLVPMVATV & 9 & HLA-A*02:01 & 2.80 \\
1OGA & GILGFVFTL & 9 & HLA-A*02:01 & 1.40 \\
3UTS & ALWGPDPAAA & 10 & HLA-A*02:01 & 2.71 \\
5C0A & MVWGPDPLYV & 10 & HLA-A*02:01 & 2.46 \\
5C0B & RQFGPDFPTI & 10 & HLA-A*02:01 & 2.03 \\
5HHO & GILEFVFTL & 9 & HLA-A*02:01 & 2.95 \\
5EU6 & YLEPGPVTV & 9 & HLA-A*02:01 & 2.02 \\
2VLR & GILGFVFTL & 9 & HLA-A*02:01 & 2.30 \\
4MJI & TAFTIPSI & 8 & HLA-B*51:01 & 2.99 \\
5NME & SLYNTVATL & 9 & HLA-A*02:01 & 2.94 \\
\midrule
1BD2 & LLFGYPVYV & 9 & HLA-A*02:01 & 2.50 \\
1LP9 & ALWGFFPVL & 9 & HLA-A*02:01 & 2.00 \\
1MI5 & FLRGRAYGL & 9 & HLA-B*08:01 & 2.50 \\
1QSE & LLFGYPRYV & 9 & HLA-A*02:01 & 2.80 \\
2AK4 & LPEPLPQGQLTAY & 13 & HLA-B*35 & 2.50 \\
2BNQ & SLLMWITQV & 9 & HLA-A*02:01 & 1.70 \\
2E7L & QLSPFPFDL & 9 & H-2L$^d$ (mouse) & 2.50 \\
2J8U & ALWGFFPVL & 9 & HLA-A*02:01 & 2.88 \\
2JCC & ALWGFFPVL & 9 & HLA-A*02:01 & 2.50 \\
2OI9 & QLSPFPFDL & 9 & H-2L$^d$ (mouse) & 2.35 \\
2PYE & SLLMWITQC & 9 & HLA-A*02:01 & 2.30 \\
2UWE & ALWGFFPVL & 9 & HLA-A*02:01 & 2.40 \\
3C60 & FEAQKAKANKAVD & 13 & H-2A (mouse, class II) & 3.05 \\
3D3V & LLFGFPVYV & 9 & HLA-A*02:01 & 2.80 \\
3H9S & MLWGYLQYV & 9 & HLA-A*02:01 & 2.70 \\
3PWP & LGYGFVNYI & 9 & HLA-A*02:01 & 2.69 \\
3QDJ & AAGIGILTV & 9 & HLA-A*02:01 & 2.30 \\
3QIB & ADLIAYLKQATKG & 13 & I-E$^k$ (mouse, class II) & 2.70 \\
4FTV & LLFGYPVYV & 9 & HLA-A*02:01 & 2.74 \\
4JFD & ELAAIGILTV & 10 & HLA-A*02:01 & 2.46 \\
4JFE & ELAGIGALTV & 10 & HLA-A*02:01 & 2.70 \\
4JFF & ELAGIGILTV & 10 & HLA-A*02:01 & 2.43 \\
4L3E & ELAGIGILTV & 10 & HLA-A*02:01 & 2.56 \\
4MNQ & ILAKFLHWL & 9 & HLA-A*02:01 & 2.74 \\
4OZG & APQPELPYPQPG & 12 & HLA-DQ (class II) & 3.00 \\
4P23 & FEAQKAKANKAVD & 13 & H-2A (mouse, class II) & 2.25 \\
4P5T & FEAQKAKANKAVD & 13 & H-2A (mouse, class II) & 3.26 \\
5E9D & ELAGIGILTV & 10 & HLA-A*02:01 & 2.51 \\
6AM5 & SMLGIGIVPV & 10 & HLA-A*02:01 & 2.39 \\
6AMU & MMWDRGLGMM & 10 & HLA-A*02:01 & 2.15 \\
\bottomrule
\end{tabular}
\caption{\textbf{The fifty-structure panel.} PDB accession, native peptide,
peptide length, MHC allele, and crystallographic resolution for every structure in the panel.
The rule separates the twenty class I HLA-A*02:01-dominated structures above from the thirty
added subsequently, which extend the panel to HLA-B, mouse H-2 and MHC class II.}
\label{tab:pdblist}
\end{table*}

\begin{table*}[p]
\renewcommand{\thetable}{S2}
\centering
\small
\begin{tabular}{@{}lcp{4.2cm}p{5.2cm}p{2.6cm}@{}}
\toprule
Model & Params & Architecture & Training data & Predicted structures \\
\midrule
ProteinMPNN & 1.66M & Message-passing GNN & PDB (\texttt{pdb\_2021aug02}), general backbones & -- \\[2pt]
ProteinMPNN (no MHC) & 1.66M & Message-passing GNN (identical to above) & Same corpus, MHC/TCR-class structures excluded & -- \\[2pt]
ESM-IF1 & 141.7M & GVP-GNN + transformer & CATH-derived backbones & +12M AF2 structures \\[2pt]
LigandMPNN & 2.62M & Message-passing GNN + ligand context & PDB, augmented with explicit ligand/heteroatom context & -- \\
\bottomrule
\end{tabular}
\caption{\textbf{Models compared.} Parameter counts verified directly from each
checkpoint's \texttt{state\_dict}. ESM-IF1 is the only tool whose training set is built substantially
from AlphaFold2-predicted structures~\citep{jumper2021alphafold} rather than solved crystal structures ($\sim$12M predicted backbones,
on top of its CATH-derived crystal set~\citep{sillitoe2021cath}).}
\label{tab:models}
\end{table*}

\begin{figure*}[p]
\renewcommand{\thefigure}{S1a}
\centering
\includegraphics[width=0.75\textwidth]{figures/paper/figS1a_peptide_length_distribution.png}
\caption{\textbf{Peptide length distribution} across all 50 structures in the
panel.}
\label{fig:supp1a}
\end{figure*}

\begin{figure*}[p]
\renewcommand{\thefigure}{S1b}
\centering
\includegraphics[width=0.75\textwidth]{figures/paper/figS1b_resolution_distribution.png}
\caption{\textbf{Crystallographic resolution per structure.}}
\label{fig:supp1b}
\end{figure*}

\begin{figure*}[p]
\renewcommand{\thefigure}{S1c}
\centering
\includegraphics[width=0.75\textwidth]{figures/paper/figS1c_mhc_contact_density.png}
\caption{\textbf{MHC residue vs.\ peptide contact density,} pooled across all 50 structures.}
\label{fig:supp1c}
\end{figure*}

\begin{figure*}[p]
\renewcommand{\thefigure}{S2}
\centering
\includegraphics[width=0.85\textwidth]{figures/paper/figS2_total_vs_unique_designs.png}
\caption{\textbf{Extreme redundancy at $T{=}0.1$.} 10,000 raw draws per
(structure, model); the unique fraction is $0.50\%$ on average and never exceeds a few percent.}
\label{fig:supp2}
\end{figure*}

\begin{figure*}[p]
\renewcommand{\thefigure}{S3}
\centering
\includegraphics[width=0.85\textwidth]{figures/paper/figS3_design_logos_pooled.png}
\caption{\textbf{Design sequence logos, pooled by peptide length, one column per
model.}}
\label{fig:supp3}
\end{figure*}

\begin{figure*}[p]
\renewcommand{\thefigure}{S4a}
\centering
\includegraphics[width=0.7\textwidth]{figures/paper/figS4a_design_logos_proteinmpnn.png}
\caption{\textbf{Per-structure design logos, ProteinMPNN.} All 50 structures,
ordered by native peptide then resolution.}
\label{fig:supp4a}
\end{figure*}

\begin{figure*}[p]
\renewcommand{\thefigure}{S4b}
\centering
\includegraphics[width=0.7\textwidth]{figures/paper/figS4b_design_logos_nomhc.png}
\caption{\textbf{Per-structure design logos, ProteinMPNN (no MHC).}}
\label{fig:supp4b}
\end{figure*}

\begin{figure*}[p]
\renewcommand{\thefigure}{S4c}
\centering
\includegraphics[width=0.7\textwidth]{figures/paper/figS4c_design_logos_esmif1.png}
\caption{\textbf{Per-structure design logos, ESM-IF1.}}
\label{fig:supp4c}
\end{figure*}

\begin{figure*}[p]
\renewcommand{\thefigure}{S4d}
\centering
\includegraphics[width=0.7\textwidth]{figures/paper/figS4d_design_logos_ligandmpnn.png}
\caption{\textbf{Per-structure design logos, LigandMPNN.}}
\label{fig:supp4d}
\end{figure*}

\begin{figure*}[p]
\renewcommand{\thefigure}{S5}
\centering
\includegraphics[width=0.75\textwidth]{figures/paper/figS5_unique_designs_per_model.png}
\caption{\textbf{Unique designs per model, mean $\pm$ SD across structures,}
pMHC+TCR vs.\ pMHC only.}
\label{fig:supp5}
\end{figure*}

\begin{figure*}[p]
\renewcommand{\thefigure}{S6}
\centering
\includegraphics[width=0.75\textwidth]{figures/paper/figS6_recovery_and_unique_vs_resolution.png}
\caption{\textbf{Resolution and unique-design count vs.\ mean recovery, pooled
across the panel.} $r=-0.32$ ($p=0.031$) and $r=0.41$ ($p=0.005$), $n=45$. Five structures are
excluded because LigandMPNN produced no designs for them, and summing unique counts over three
models rather than four would understate their diversity.}
\label{fig:supp6}
\end{figure*}

\begin{figure*}[p]
\renewcommand{\thefigure}{S7}
\centering
\includegraphics[width=0.75\textwidth]{figures/paper/figS7_replicate_groups_vs_quality.png}
\caption{\textbf{Resolution dependence in a controlled, same-peptide replicate
set.} Four structures (2P5W, 2BNR, 2F53, 2F54), identical native peptide, near-identical TCR clone; every
model plotted separately.}
\label{fig:supp7}
\end{figure*}

\begin{figure*}[p]
\renewcommand{\thefigure}{S8a}
\centering
\includegraphics[width=0.75\textwidth]{figures/paper/figS8a_esmcba_vs_mhcflurry_no_tcr.png}
\caption{\textbf{ESMCBA vs.\ MHCflurry agreement, pMHC-only designs,}
by model.}
\label{fig:supp8a}
\end{figure*}

\begin{figure*}[p]
\renewcommand{\thefigure}{S8b}
\centering
\includegraphics[width=0.75\textwidth]{figures/paper/figS8b_resolution_by_model.png}
\caption{\textbf{Per-model resolution/design-count relationship,} full panel.}
\label{fig:supp8b}
\end{figure*}

\begin{figure*}[p]
\renewcommand{\thefigure}{S9a}
\centering
\includegraphics[width=0.85\textwidth]{figures/paper/figS9a_own_score_shift_pooled.png}
\caption{\textbf{Each model's own confidence score for its designs shifts when
the TCR is removed,} pooled across all forty HLA-A*02:01 structures. ProteinMPNN and ProteinMPNN (no
MHC): average per-residue negative log-likelihood, lower is more confident. LigandMPNN: overall
confidence, higher is more confident. All three shift toward lower self-assessed confidence without the
TCR (Mann-Whitney $p<10^{-200}$ in all three cases).}
\label{fig:supp9a}
\end{figure*}

\begin{figure*}[p]
\renewcommand{\thefigure}{S9b}
\centering
\includegraphics[width=0.7\textwidth]{figures/paper/figS9b_own_score_shift_by_structure.png}
\caption{\textbf{The own-score shift is not uniform across the panel.}
Bonferroni-corrected significance (full vs.\ mhconly) per structure per model; dark cells are structures
where the model's own-score distribution did not shift.}
\label{fig:supp9b}
\end{figure*}

\begin{figure*}[p]
\renewcommand{\thefigure}{S10}
\centering
\includegraphics[width=0.75\textwidth]{figures/paper/figS10_strong_binder_rate_by_model.png}
\caption{\textbf{Removing the TCR raises predicted affinity for every model, by very different
amounts.} Strong-binder rate (MHCflurry IC50 $<$ 500\,nM) per model in both structural arms; solid
bars are pMHC+TCR, hatched bars pMHC only. Brackets give a Fisher's exact test of that model's own
shift between arms, Bonferroni-corrected across the four models. Every model shifts up significantly,
but the magnitude separates them: ProteinMPNN (no MHC) gains $25.7$ percentage points and LigandMPNN
$10.5$, against $5.3$ for ProteinMPNN and $4.0$ for ESM-IF1. ESM-IF1 is the lowest-affinity model in
both arms. The model ranking is preserved between arms, so the omnibus difference between models
(reported above each arm) is not an artifact of which context is used.}
\label{fig:supp10}
\end{figure*}

\begin{figure*}[p]
\renewcommand{\thefigure}{S11}
\centering
\includegraphics[width=\textwidth]{figures/paper/figS11_umap_six_colorings.png}
\caption{\textbf{The sampled sequence space, colored by predicted binding
affinity.} One UMAP (ESMCBA embeddings, fit jointly over every unique design), recolored six ways. All
four models occupy the same overlapping landscape rather than four separate clusters, and MHCflurry and
ESMCBA, built differently and trained differently, trace the identical gradient across it.}
\label{fig:supp11}
\end{figure*}

\begin{figure*}[p]
\renewcommand{\thefigure}{S12}
\centering
\includegraphics[width=0.62\textwidth]{figures/paper/figS12_native_anchor_hydrophobicity.png}
\caption{\textbf{Both MHC anchors are more hydrophobic than the peptide interior they sit alongside.}
Kyte-Doolittle hydrophobicity of the native residue at P2 and P$\Omega$ for each of the 40 HLA-A*02:01
structures, plotted against that same peptide's mean over its own interior positions (P3 through the
second-to-last, excluding P1). Grey lines connect the three values belonging to one peptide; anchor
points are colored by side-chain class. Both anchors sit well above their own interior
(paired Wilcoxon, $p<0.0001$ for both, $n=40$), so the two pockets do not
differ in whether they prefer hydrophobic residues. They differ in how strictly the designs honor that
preference, which is the distinction reported in Section~\ref{sec:recovery}. The points at
P$\Omega$ colored as polar are chiefly the cysteine-terminated SLLMWITQC/NY-ESO-1 replicates, a residue whose
discrete classification understates its hydrophobicity.}
\label{fig:supp12}
\end{figure*}

\begin{figure*}[p]
\renewcommand{\thefigure}{S13}
\centering
\includegraphics[width=\textwidth]{figures/paper/figS13_per_position_recovery_10mer.png}
\caption{\textbf{Per-position recovery, decamers.} As Figure~\ref{fig:fig1}, for the 11 crystals with a
10-residue epitope. The P2 anchor again recovers above the interior, and the C-terminal anchor P$\Omega$
falls to $0.08$, lower than in the nonamers.}
\label{fig:supp_len10}
\end{figure*}

\begin{figure*}[p]
\renewcommand{\thefigure}{S14}
\centering
\includegraphics[width=\textwidth]{figures/paper/figS14_per_position_recovery_13mer.png}
\caption{\textbf{Per-position recovery, 13-mers.} As Figure~\ref{fig:fig1}, for the 5 crystals with a
13-residue epitope. Four of the five are MHC class II complexes, which have no P2 or P$\Omega$ anchor
pocket, and the anchor structure seen at class I is correspondingly absent.}
\label{fig:supp_len13}
\end{figure*}

\begin{figure*}[p]
\renewcommand{\thefigure}{S15}
\centering
\includegraphics[width=\textwidth]{figures/paper/figS15_skempi_ddg_vs_recovery.png}
\caption{\textbf{Measured binding effect vs.\ recovery, pooled over every SKEMPI complex with designs.}
Each point is one (complex, peptide position) with a measured single-mutation
$\Delta\Delta G$; positive $\Delta\Delta G$ means the mutation weakens binding, marking a position the
interface depends on. Left, recovery with the TCR present. Middle, the recovery gain from TCR context.
Right, that gain against $\Delta\Delta G$ directly. 55 positions across 17 complexes, against the five
positions in one crystal available before the panel was extended. $\Delta\Delta G$ marks positions the
interface depends on, which includes MHC anchors as well as TCR contacts.}
\label{fig:supp_skempi_pooled}
\end{figure*}


\begin{figure*}[p]
\renewcommand{\thefigure}{S16}
\centering
\includegraphics[width=\textwidth]{figures/paper/figS16_dmf5_score_vs_kd.png}
\caption{\textbf{The score carries no information about measured affinity.} LigandMPNN score against
published $K_D$ for the nine DMF5 peptides with an actual measured value, scored on 6AM5 (left) and
6AMU (right); points are colored by melting temperature and each structure's own crystallized peptide
is circled. Four further peptides are excluded because their $K_D$ is reported only as $>500\,\mu$M or
as not detected, which is an assay ceiling rather than a measurement. Neither backbone shows a
relationship. Contrast this with Figure~\ref{fig:fig4}, where the same scores
separate the two peptide families cleanly: these models place a peptide, they do not rank it.}
\label{fig:supp_6am_kd}
\end{figure*}

\begin{figure*}[p]
\renewcommand{\thefigure}{S17}
\centering
\includegraphics[width=\textwidth]{figures/paper/figS17_anchor_non_redundant_sets.png}
\caption{\textbf{The anchor split is not an artefact of the panel's redundancy.} The panel holds 50
structures but only 33 distinct peptides and 26 distinct CDR3$\beta$ sequences, so a pooled statistic
weights an epitope by how often it was crystallised. Left, P2, interior and P$\Omega$ recovery on
progressively stricter sets: human MHC class I only, then one structure per TCR family (CDR3$\beta$
clustered so clonotypes differing by one or two substitutions collapse), then one per peptide, then
restricted to HLA-A*02:01; each group is represented by its best-resolution member, chosen
independently of any score. Right, the same result as the two gaps that define it. Both survive, and
the P$\Omega$ deficit deepens from $0.23$ to $0.35$ as repeated epitopes are removed. The one
exception is P2 on the TCR-only set, which at 17 structures is underpowered rather than contradictory.}
\label{fig:supp_dedup}
\end{figure*}

\end{document}
